% Preamble needs: \usepackage{xcolor} \usepackage{soul}
% Yellow (\hl, \colorbox{yellow}) = new; red strikethrough (\textcolor{red}{\st{...}}) = removed.
\subsection{Production}
The production sector in country $X$ consists of three \textcolor{red}{\st{representative firms that produce, aggregate consumption goods and a capital producer.}} \hl{layers of firms and a capital producer. Competitive intermediate producers combine capital and hours into a homogeneous intermediate good; retailers differentiate it; and a competitive aggregator bundles the varieties into country $X$'s home good, the good each country trades. Households combine the two home goods into a consumption basket; investment and government purchases use the home good only.}
\paragraph{Consumption Baskets and the terms of trade:}
The consumption basket of country $D$ aggregates the home and foreign goods with an
Armington/CES aggregator:
\begin{equation}
C_{D,t}=\Bigl[\varpi_D^{1/\eta}c_{DD,t}^{\frac{\eta-1}{\eta}}
+(1-\varpi_D)^{1/\eta}c_{FD,t}^{\frac{\eta-1}{\eta}}\Bigr]^{\frac{\eta}{\eta-1}},
\label{eq:ces-aggregator}
\end{equation}
with home bias $\varpi_D$ and trade elasticity \textcolor{red}{\st{$\eta_D$}}\hl{$\eta$}. $C_{F,t}$ is defined
symmetrically, with weight $\varpi_F$ on $c_{FF,t}$. Normalising the price of the
$D$-good to one and letting $p_t$ be the relative price of the $F$-good, cost
minimisation gives the price index and the import demand schedule
\begin{align}
P^{c}_{D,t}&=\Bigl[\varpi_D+(1-\varpi_D)\,p_t^{\,1-\eta}\Bigr]^{\frac{1}{1-\eta}},
&
IM_{D,t}&=(1-\varpi_D)\Bigl(\frac{P^{c}_{D,t}}{p_t}\Bigr)^{\eta}C_{D,t},
\label{eq:ces-D}\\[3pt]
P^{c}_{F,t}&=\Bigl[\varpi_F+(1-\varpi_F)\,p_t^{\,\eta-1}\Bigr]^{\frac{1}{1-\eta}},
&
IM_{F,t}&=(1-\varpi_F)\bigl(P^{c}_{F,t}\,p_t\bigr)^{\colorbox{yellow}{$\scriptstyle\eta$}}C_{F,t},
\label{eq:ces-F}
\end{align}
where $P^{c}_{F,t}$ is expressed in $F$-good units and the $D$-good costs $1/p_t$
there. With unequal masses, the two home-goods weights satisfy:
\begin{equation}
1-\varpi_F=\frac{1-\varpi_D}{\varkappa},
\label{eq:home-bias-size}
\end{equation}
and net exports, in own-good units per capita, are:
\begin{equation}
NX_{D,t}=\varkappa\,IM_{F,t}-p_t\,IM_{D,t},
\qquad
NX_{F,t}=\frac{IM_{D,t}}{\varkappa}-\frac{IM_{F,t}}{p_t}.
\label{eq:nx}
\end{equation}

\paragraph{Retailers and marginal cost:}
A unit continuum of retailers buys the homogeneous intermediate good, differentiates it costlessly, and sells under monopolistic competition to a \textcolor{red}{\st{CES final-goods aggregator firm}} \hl{competitive aggregator that bundles the varieties into the home good,} \colorbox{yellow}{$Y_{X,t}=\bigl[\int_0^1 y_{jX,t}^{(\epsilon_X-1)/\epsilon_X}\,\mathrm{d}j\bigr]^{\epsilon_X/(\epsilon_X-1)}$}\hl{,} with elasticity $\epsilon_X$. \hl{A retailer's marginal cost is the price of the intermediate good, $mc_X$, in units of the home good.} Prices are fully flexible, so every retailer charges the constant markup $\epsilon_X/(\epsilon_X-1)$ over marginal cost\hl{; in the symmetric equilibrium every retailer's price equals that of the home good, one,} and real marginal cost is constant\textcolor{red}{\st{;}}\hl{,}
\begin{equation}
\colorbox{yellow}{$\displaystyle mc_X=\frac{\epsilon_X-1}{\epsilon_X}$}
\label{eq:mc}
\end{equation}
\hl{With no price dispersion, home-good output equals intermediate output $Y_{X,t}$ one for one.} Retail profits $(1-mc_X)Y_{X,t}$ are rebated lump-sum to households.

\paragraph{Intermediate producers and the working-capital wedge:}
\label{sec:wc}

The intermediate producer operates a Cobb--Douglas technology in the predetermined capital stock ($K_{X,t}$) and hours ($N_{X,t}$):
\begin{equation}
Y_{X,t}=Z_{X,t}\,K_{X,t}^{\alpha_X}\,N_{X,t}^{1-\alpha_X},
\label{eq:production}
\end{equation}
where $K_{X,t}$ was purchased by banks at $t-1$ \hl{and is rented from them at $mpk_{X,t}$; output is sold to retailers at $mc_X$}. The firm must pre-finance a fraction $\zeta_X$ of its wage bill with intra-period bank credit at the gross rate $1+r^{wc}_{X,t}$ \citep{neumeyer2005business}. Thus the static problem is of the firm:
\begin{equation}
\max_{N_{X,t}\colorbox{yellow}{$\scriptstyle ,\,K_{X,t}$}}\;\;
mc_X\,Z_{X,t}K_{X,t}^{\alpha_X}N_{X,t}^{1-\alpha_X}
\;-\;\bigl(1+\zeta_X r^{wc}_{X,t}\bigr)\,w_{X,t}N_{X,t}
\;-\;mpk_{X,t}K_{X,t},
\label{eq:firm-problem}
\end{equation}
delivering the factor-demand conditions:
\begin{align}
w_{X,t}&=\frac{mc_X(1-\alpha_X)\,Y_{X,t}/N_{X,t}}{1+\zeta_X\,r^{wc}_{X,t}},
\label{eq:labour-demand}\\[3pt]
mpk_{X,t}&=mc_X\,\alpha_X\,\frac{Y_{X,t}}{K_{X,t}},
\label{eq:mpk}
\end{align}
and a loan demand of:
\begin{equation}
L_{X,t}=\zeta_X\,w_{X,t}N_{X,t}.
\label{eq:loan-demand}
\end{equation}
\hl{With constant returns the producer earns zero profit: the financed wage bill absorbs $mc_X(1-\alpha_X)Y_{X,t}$ and the rental $mc_X\alpha_X Y_{X,t}$ accrues to the banks that own the capital.}
Notice that equation~\eqref{eq:labour-demand} is the \textbf{only} channel through which
financial spreads, affecting $r_{X,t}^{wc}$, reach output on impact.

\paragraph{Capital producers, Tobin's \textcolor{red}{\st{q}} \colorbox{yellow}{$Q$}, and the return on capital:}
Capital producers convert the home good into installed capital, subject to adjustment costs \cite{jermann1998asset} and depreciation rate $\delta_X$. With $\iota_{X,t}\equiv I_{X,t}/K_{X,t}$, the law of motion is:
\begin{equation}
K_{X,t+1}=(1-\delta_X)K_{X,t}+\Phi_X(\iota_{X,t})\,K_{X,t},
\qquad
\Phi_X(\iota)=\gamma_{0,X}\,\iota^{\,1-\xi_X}+\gamma_{1,X}.
\label{eq:capital-lom}
\end{equation}
The capital producer solves $\max_{\iota}\;Q_{X,t}\Phi_X(\iota)K_{X,t}-\iota K_{X,t}$, whose first-order condition delivers marginal Tobin's $Q$, in home goods units, and, inverting
\eqref{eq:capital-lom}, the investment rate:
\begin{equation}
Q_{X,t}=\frac{1}{\Phi_X'(\iota_{X,t})}
      =\frac{\iota_{X,t}^{\;\xi_X}}{\gamma_{0,X}(1-\xi_X)},
\qquad
\iota_{X,t}=\left[\frac{K_{X,t+1}/K_{X,t}-(1-\delta_X)-\gamma_{1,X}}
                       {\gamma_{0,X}}\right]^{\frac{1}{1-\xi_X}} .
\label{eq:tobins-q}
\end{equation}
Capital producers' rents, $\Pi^{k}_{X,t}=Q_{X,t}\bigl[K_{X,t+1}-(1-\delta_X)K_{X,t}\bigr]-I_{X,t}$, are rebated to households. The parameter $\xi_X$ is the elasticity of $Q$ with respect to $I/K$. The bank holds the capital, so the realised gross return between $t$ and $t+1$ on a unit of capital purchased at $t$ is:
\begin{equation}
\;1+r^{k}_{X,t+1}=\frac{mpk_{X,t+1}+(1-\delta_X)\,Q_{X,t+1}}{Q_{X,t}}\
\label{eq:rk}
\end{equation}
